\section{Методы наблюдения и моделирования}

\subsection{Наземные инструменты}

Наземные наблюдения включают авроральные камеры, фотометры, радары некогерентного рассеяния (EISCAT) и магнитометры. Сети камер, такие как THEMIS Ground Network, отслеживают движение сияний в реальном времени.

\subsection{Космические миссии}

Космические аппараты позволяют измерять процессы в магнитосфере:
\begin{itemize}
    \item \textbf{THEMIS} (NASA, 2007) --- пять спутников, локализовавших точку старта суббури~\cite{angelopoulos2008}.
    \item \textbf{Cluster} (ESA, 2000) --- четыре спутника для изучения трёхмерных структур.
    \item \textbf{MMS} (NASA, 2015) --- высокое временное разрешение для изучения пересоединения~\cite{burch2016}.
\end{itemize}

\subsection{Программная реализация моделирования движения частиц}

Ниже представлен фрагмент программного кода на Python, моделирующий движение заряженной частицы в магнитном поле Земли методом Бориса.

\begin{minted}[frame=lines, framesep=2mm, fontsize=\small, linenos]{python}
import numpy as np
import matplotlib.pyplot as plt

"""
Моделирование движения заряженной частицы в магнитном поле Земли
(метод Бориса для уравнения движения в магнитном поле)
"""

def boris_step(r, v, B, q, m, dt):
    """
    Один шаг метода Бориса для движения частицы в магнитном поле.
    
    Параметры:
        r - позиция (x, y, z)
        v - скорость (vx, vy, vz)
        B - магнитное поле в точке r
        q - заряд частицы
        m - масса частицы
        dt - шаг по времени
    """
    # Полу-шаг для скорости
    v_minus = v + 0.5 * dt * (q / m) * np.cross(v, B)
    
    # Поворот в магнитном поле
    t = (q * dt * 0.5 / m) * B
    s = 2 * t / (1 + np.dot(t, t))
    v_prime = v_minus + np.cross(v_minus, t)
    v_plus = v_minus + np.cross(v_prime, s)
    
    # Полный шаг
    v_new = v_plus + 0.5 * dt * (q / m) * np.cross(v_plus, B)
    r_new = r + v_new * dt
    return r_new, v_new

# Параметры частицы (электрон)
q = -1.602e-19      # заряд, Кл
m = 9.109e-31       # масса, кг
v0 = np.array([1e6, 0, 0])   # начальная скорость, м/с
r0 = np.array([6.4e6, 0, 0]) # начальная позиция, м
dt = 1e-5           # шаг, с
n_steps = 10000     # количество шагов

# Массивы для траектории
r_traj = [r0.copy()]
v_traj = [v0.copy()]

# Основной цикл
r = r0.copy()
v = v0.copy()
for _ in range(n_steps):
    # Простейшая модель дипольного поля
    R = np.linalg.norm(r)
    B = np.array([0, 0, 1e-5 * (6.4e6 / R)**3])  # Bz только для примера
    r, v = boris_step(r, v, B, q, m, dt)
    r_traj.append(r.copy())
    v_traj.append(v.copy())

print(f"Моделирование завершено. Количество шагов: {n_steps}")
\end{minted}

\subsection{Математическое моделирование}

Три основных подхода: глобальные MHD-модели (LFM, BATSRUS), кинетические модели (трассировка частиц) и гибридные модели (MHD + кинетика). Прогнозы авроральной активности достигают около 70\% точности.